
\subsection{算子 $\mathfrak q$ 的构造 IV：
定理 3.8.32 的证明收束}
\label{subsec:7-4-5}
\liftIndex{I359}

为把第 7.4.2--7.4.4 小节的结果
应用于定理~\ref{thm:3-8-32} 的证明，
需要第 7.4.3、7.4.4 小节中结果的
$A_{n,K}$ 版本。

设
$G\subset\mathbb R_{\ge0}\times2\mathbb Z$
是满足条件 3.1.6 的幺半群。
对 $\beta\in G$，
如定义~\ref{def:7-2-61} 定义
$\|\beta\|$，
并如定义~\ref{def:7-2-63}
在 $G\times\mathbb Z_{\ge0}$ 上定义关系 $>$。
设 $(C,\mathfrak m)$
为一个 $G$-有隙滤过 $A_{n,K}$ 代数。

\begin{definition}{7.4.107}\label{def:7-4-107}
设
\[
(\varphi_k)_{k=0,1,\ldots}
\in
\prod_{k=0}^{\infty}
\operatorname{Hom}\bigl(B_k(C[1]),C[1]\bigr).
\]
若每个 $\varphi_k$ 都可分解为
\[
\varphi_k
=\sum_{\beta\in G}
T^{E(\beta)}e^{\mu(\beta)/2}\varphi_{k,\beta},
\]
其中
\[
\varphi_{k,\beta}\in
\operatorname{Hom}\bigl(
B_k(\overline C[1]),
\overline C[1]
\bigr)
\]
是 $\mathbb Q$-线性的，
则称该元素为 $G$-有隙的。
记所有 $G$-有隙元素所成集合为
\[
CH_G(C,C)
\bigl(:=\operatorname{Der}B_G(C[1])\bigr).
\]

令
$\mathfrak I_{n',K'}\subset CH_G(C,C)$
为所有满足下述条件的
$(\varphi_k)_{k=0,1,\ldots}$ 所成子集：
\[
\varphi_{k,\beta}=0
\qquad
\text{若 }(\beta,k)\precsim(n',K').
\]
含幺版本
\[
CH_G^{\mathrm{unit}}(C,C)
=\operatorname{Der}B_G^{\mathrm{unit}}(C[1]),
\qquad
\mathfrak I_{n',K'}^{\mathrm{unit}}
\]
也以同样方式定义。
\end{definition}

\begin{lemma}{7.4.108}\label{lem:7-4-108}
设 $(C,\mathfrak m)$
为一个 $G$-有隙滤过 $A_{n,K}$ 代数。

\par\noindent\textup{(1)}
$\operatorname{Der}B_G(C[1])$
是 $\operatorname{Der}B(C[1])$
的分次李子代数。

\par\noindent\textup{(2)}
$\mathfrak I_{n,K}$
是 $\operatorname{Der}B_G(C[1])$ 的理想。

\par\noindent\textup{(3)}
\[
\operatorname{Im}(\delta\circ\delta)
\subseteq\mathfrak I_{n,K},
\]
其中 $\delta$ 由式~\eqref{eq:7-4-27} 定义。

\par\noindent\textup{(4)}
若 $(n',K')\precsim(n,K)$，则
\[
\frac{\operatorname{Der}B_G(C[1])}
{\mathfrak I_{n',K'}}
\]
是微分分次李代数。

含幺版本的 (1)、(2) 也成立。
\end{lemma}

\begin{proof}
(1) 来自 $G$ 是幺半群这一事实。
(2) 由引理~\ref{lem:7-2-65} 得到。
(3)、(4) 的证明是容易的。
\end{proof}

\begin{definition}{7.4.109}\label{def:7-4-109}
置
\[
\operatorname{Der}B_{G;n,K}(C[1])
=
\frac{\operatorname{Der}B_G(C[1])}
{\mathfrak I_{n,K}}.
\]
\end{definition}
由引理~\ref{lem:7-4-108}，
这是一个微分分次李代数；
这里取 $(n',K')=(n,K)$。

下文省略 $G$，
把 $G$-有隙滤过 $A_{n,K}$ 代数简称为
滤过 $A_{n,K}$ 代数，
其他名称亦同。
也从
\[
\frac{\operatorname{Der}B_G(C[1])}
{\mathfrak I_{n,K}}
\]
中省略 $G$，
简写为 $\operatorname{Der}B_{n,K}(C[1])$。

下面是定理~\ref{thm:7-4-41}
的 $(n,K)$ 版本。

\begin{theorem}{7.4.41$'$}\label{thm:7-4-41-prime}
若 $(C,\mathfrak m)$ 与
$(C',\mathfrak m')$\translatornote{源书靶侧写作 $(C',\mathfrak m)$，结构记号漏撇；本定理是定理 7.4.41 的 $(n,K)$ 版本，故补为 $\mathfrak m'$。见疑误 A-0580。}
是 $A_{n,K}$-同伦等价的，
则微分分次李代数
$\operatorname{Der}B_{n,K}(C[1])$
与
$\operatorname{Der}B_{n,K}(C'[1])$
严格滤过地 $L_\infty$ 同伦等价。
\end{theorem}

证明与定理~\ref{thm:7-4-41} 相同，故略。

下面用上述记号，
把引理~\ref{lem:7-4-17}
改写成 $L_l$ 结构的语言。
取数列 $\mathcal K_\lambda>0$、
$\delta_\lambda>0$，满足
\[
\mathcal K_\lambda\nearrow\infty,
\qquad
\delta_\lambda\searrow0
\qquad(\lambda\to\infty).
\]
将引理~\ref{lem:7-4-17}
应用于
$\delta=\delta_\lambda$、
$K=\mathcal K_\lambda$。
于是对
\[
(\beta,k)\precsim
(\mathcal K_\lambda,\mathcal K_\lambda),
\qquad
\ell<\mathcal K_\lambda,
\]
定义下面的算子
$\mathfrak q^{\mathrm{geo}}_{\ell,k,\beta}$。

令
$P_i\in\mathcal X_{1+,(\mathcal K_\lambda)}(L)$、
$Q_i\in\mathcal X(M)$\translatornote{源书把前一个集合写作未定义的 $\mathcal X_{+,(\mathcal K_\lambda)}(L)$；本段正在应用引理 7.4.17，而该引理定义的是 $\mathcal X_{1+,(\mathcal K_\lambda)}(L)$，故补下标 $1$。见疑误 A-0581。}。
置
\begin{equation}
\begin{aligned}
\mathfrak q^{\mathrm{geo}}_{\ell,k,\beta}
(\vec Q)(\vec P)
=\operatorname{ev}_{0*}\Bigl(
\mathcal M^{\mathfrak d}_{k+1,\ell}
(L;\beta;\vec Q,\vec P)^{
\mathfrak s_{\mathfrak d,k,\beta,\vec Q,\vec P}}
\Bigr).
\end{aligned}
\tag{7.4.110}\label{eq:7-4-110}
\end{equation}
注意，式~\eqref{eq:7-4-110}
中的映射只得到部分定义。

接着仿照使用定理~\ref{thm:5-4-2}
（的证明）的方式，
使用定理~\ref{thm:7-4-41}
（的证明）并应用第~\ref{subsec:7-2-7} 小节的论证，
得到
\begin{equation}
\mathfrak q_{\ell,k,\beta}:
E_\ell\bigl(C(M;\Lambda_{0,\mathrm{nov}})[2]\bigr)
\otimes B_k\bigl(C(L;\lambda)[1]\bigr)
\longrightarrow C(L;\lambda)[1]
\tag{7.4.111}\label{eq:7-4-111}
\end{equation}
其中
\[
\ell\precsim l(\lambda),
\qquad
(\beta,k)\precsim(n(\lambda),K(\lambda)).
\]
这里
\[
l(\lambda)\nearrow\infty,
\qquad
n(\lambda)\nearrow\infty,
\qquad
K(\lambda)\nearrow\infty,
\]
而 $C(L;\lambda)$
是由
$\mathcal X_{1+,(g(\lambda))}(L)$
生成的自由 $\Lambda_{0,\mathrm{nov}}$-模，
其中 $g(\lambda)\nearrow\infty$；
源书此处同样把该生成元集合写成
$\mathcal X_{+,(g(\lambda))}(L)$，
这里依 A-0581 补下标 $1$。

算子 $\mathfrak q_{\ell,k,\beta}$
具有引理~\ref{lem:7-4-113} 所述的性质。
定义
\[
\widehat{\mathfrak q}^{\circ}_\lambda:
\bigoplus_{\ell\precsim l(\lambda)}
E_\ell\bigl(C(M;\Lambda_{0,\mathrm{nov}})[2]\bigr)
\longrightarrow
\operatorname{Der}B_{n(\lambda),K(\lambda)}
\bigl(C(L;\lambda)[1]\bigr)
\]
如下：
\begin{equation}
\widehat{\mathfrak q}^{\circ}_{\lambda;\ell}(\vec Q)
\equiv
\widehat{\mathfrak q}_\ell
\bigl(\operatorname{Op}(\vec Q)\bigr)
\pmod{\mathfrak I_{n(\lambda),K(\lambda)}}.
\tag{7.4.112}\label{eq:7-4-112}
\end{equation}

\begin{lemma}{7.4.113}\label{lem:7-4-113}
$\widehat{\mathfrak q}^{\circ}_\lambda$
是严格滤过 $L_{l(\lambda)}$ 同态。
\end{lemma}

\begin{remark}{7.4.114}\label{rem:7-4-114}
注意，对严格滤过 $L_\infty$ 代数，
可以定义严格滤过 $L_{l(\lambda)}$ 同态的概念，
因为在严格情形下数目滤过得以保持。
\end{remark}

我们将采用与第~\ref{subsec:7-2-8} 小节
类似的论证，
并应用引理~\ref{lem:7-4-113}，
以得到滤过 $L_\infty$ 同态
\[
\mathfrak q^{\circ}:
C(M;\Lambda_{0,\mathrm{nov}})[1]
\longrightarrow
\operatorname{Der}B_G\bigl(C(L;\lambda)[1]\bigr).
\]
为此，需要证明各个 $\lambda$ 所对应的
$\widehat{\mathfrak q}^{\circ}_\lambda$
彼此相容。
我们将应用命题~\ref{prop:4-6-55}
证明这种相容性；
该命题中出现的
$\mathcal X_{3,+}(L')$ 的构造
与第~\ref{subsec:7-2-9} 小节相似。
为将命题~\ref{prop:4-6-55}
用于这里，
先须以类似第~\ref{subsec:5-4-5} 小节的方式
改述引理~\ref{lem:7-4-113}。

取
$\mathcal H^{\mathrm{amb}}
\subset C(M;\Lambda_{0,\mathrm{nov}})$
为由表示同调群
$H(M;\Lambda_{0,\mathrm{nov}})$
的闭链组成的子模。
选取
\[
\{v_1,\ldots,v_b\}
\]
作为
\[
\mathcal H^{\mathrm{amb}}[1]_{\mathrm{odd}}
=(\mathcal H^{\mathrm{amb}})_{\mathrm{even}}
\]
的一组基，
并选取
\[
\{w_1,\ldots,w_{b'}\}
\]
作为
$\mathcal H^{\mathrm{amb}}[1]_{\mathrm{even}}$
的一组基。
把相应的形式坐标分别记为
$r_1,\ldots,r_b$
和 $s_1,\ldots,s_{b'}$。
置
\[
\mathbf b
=\sum_i r_iv_i+\sum_i s_iw_i.
\]

对
$x_i\in B(C(L;\lambda)[1])$，定义
\begin{equation}
\mathfrak m^{\mathbf b}(x_1,\ldots,x_k)
=
\sum_{\ell=0}^{l(\lambda)}
\mathfrak q_{\ell,k}
(\mathbf b,\ldots,\mathbf b)
(x_1,\ldots,x_k).
\tag{7.4.115}\label{eq:7-4-115}
\end{equation}
这里把式~\eqref{eq:7-4-115} 的右端
视为
\[
C(L;\lambda)_{\mathcal R^{\mathrm{sup}}_{l(\lambda)}}
:=
C(L;\lambda)\otimes_{\mathbb Q}
\mathcal R^{\mathrm{sup}}_{l(\lambda)}
\]
中的元素，其中
\[
\mathcal R^{\mathrm{sup}}
=
\mathbb Q[[r_1,\ldots,r_b]]
\otimes\Omega(s_1,\ldots,s_{b'})
\]
且
\[
\mathcal R^{\mathrm{sup}}_{l(\lambda)}
=
\frac{\mathcal R^{\mathrm{sup}}}
{\mathfrak{IR}_{l(\lambda)}}.
\]
这里
$\mathfrak{IR}_{l(\lambda)}$
是由次数大于 $l(\lambda)$ 的元素组成的理想；
“sup”代表 super（超）。

显然可以定义系数环为分次环
$\mathcal R^{\mathrm{sup}}_{l(\lambda)}$
的滤过 $A_{n(\lambda),K(\lambda)}$ 代数。

\begin{lemma}{7.4.116}\label{lem:7-4-116}
若
$\widehat{\mathfrak q}^{\circ}_\lambda$
是严格滤过 $L_{l(\lambda)}$ 同态，
则 $\mathfrak m^{\mathbf b}$
在
$C(L;\lambda)_{\mathcal R^{\mathrm{sup}}_{l(\lambda)}}$
上定义一个滤过
$A_{n(\lambda),K(\lambda)}$ 代数结构。
反过来，若
$C(L;\lambda)_{\mathcal R^{\mathrm{sup}}_{l(\lambda)}}$
上有一个滤过
$A_{n(\lambda),K(\lambda)}$ 代数结构，
且它在 $C(L;\lambda)$ 上诱导 $\mathfrak m$，
则它由某个严格滤过
$L_{l(\lambda)}$ 同态
$\widehat{\mathfrak q}^{\circ}_\lambda$ 诱导。
换言之，它与式~\eqref{eq:7-4-115}
中的 $\mathfrak m^{\mathbf b}$ 一致。
\end{lemma}

\begin{proof}
前一半已经证明。
为证明逆命题，
设 $\widetilde{\mathfrak m}$
是
$C(L;\lambda)_{\mathcal R^{\mathrm{sup}}_{l(\lambda)}}$
上的滤过
$A_{n(\lambda),K(\lambda)}$ 代数结构。
写成
\[
\begin{aligned}
\widetilde{\mathfrak m}(x_1,\ldots,x_n)
=\sum
&r_1^{l_1}\cdots r_b^{l_b}
s_1^{\epsilon_1}\cdots s_{b'}^{\epsilon_{b'}}\\
&\cdot
\widetilde{\mathfrak m}_
{l_1,\ldots,l_b,\epsilon_1,\ldots,\epsilon_{b'}}
(x_1,\ldots,x_n),
\end{aligned}
\]
其中 $l_i=0,1,2,\ldots$，
$\epsilon_i=0$ 或 $1$。
定义
\[
\begin{aligned}
&\mathfrak q_{\ell,n}
\bigl(
v_1^{\otimes l_1},\ldots,v_b^{\otimes l_b},
w_1^{\otimes\epsilon_1},\ldots,
w_{b'}^{\otimes\epsilon_{b'}}
\bigr)(x_1,\ldots,x_n)\\
&\qquad=
\frac{l_1!\cdots l_b!}
{\bigl(\sum_i l_i\bigr)!}
\widetilde{\mathfrak m}_
{l_1,\ldots,l_b,\epsilon_1,\ldots,\epsilon_{b'}}
(x_1,\ldots,x_n),
\end{aligned}
\]
其中
\[
\ell=\sum_i l_i+\sum_i\epsilon_i.
\]
容易看出，
由这个 $\mathfrak q_{\ell,n}$
经式~\eqref{eq:7-4-115}
定义的 $\mathfrak m^{\mathbf b}$
就是 $\widetilde{\mathfrak m}$\translatornote{源书总结句把刚定义的 $\mathfrak q_{\ell,n}$ 改写为 $\mathfrak q_{\ell,m}$；本段始终以 $n$ 表示输入数，故沿用显示式的下标。见疑误 A-0584。}。
\end{proof}

\begin{lemma}{7.4.117}\label{lem:7-4-117}
对 $\lambda,\lambda'$，存在
\[
n'(\lambda,\lambda')\nearrow\infty,
\qquad
K'(\lambda,\lambda')\nearrow\infty,
\qquad
l'(\lambda,\lambda')\nearrow\infty,
\]
使
\[
\frac{
C(L;\lambda)_{\mathcal R^{\mathrm{sup}}_{l(\lambda)}}
}{
\mathfrak{IR}_{l'(\lambda,\lambda')}
}
\quad
\left(
\cong
C(L;\lambda)_
{\mathcal R^{\mathrm{sup}}_{l'(\lambda,\lambda')}}
\right)
\]
与
\[
\frac{
C(L';\lambda')_{\mathcal R^{\mathrm{sup}}_{l(\lambda')}}
}{
\mathfrak{IR}_{l'(\lambda,\lambda')}
}
\quad
\left(
\cong
C(L';\lambda')_
{\mathcal R^{\mathrm{sup}}_{l'(\lambda,\lambda')}}
\right)
\]
是
$A_{n'(\lambda,\lambda'),K'(\lambda,\lambda')}$
同伦等价的。
\end{lemma}

\begin{proof}
仿照定理~\ref{thm:4-6-52} 的证明，
使用命题~\ref{prop:4-6-55} 的
$A_{n(\lambda),K(\lambda)}$ 版本，
即可构造同伦等价。
细节留待读者完成。
\end{proof}

现在可以完成定理~\ref{thm:3-8-32} 的证明。
取 $\lambda_i\nearrow\infty$（$i\to\infty$）。
不失一般性，可假定
\[
\begin{aligned}
n'(\lambda_i,\lambda_{i+1})&=n(\lambda_i),\\
K'(\lambda_i,\lambda_{i+1})&=K(\lambda_i),\\
l'(\lambda_i,\lambda_{i+1})&=l(\lambda_i).
\end{aligned}
\]
由引理~\ref{lem:7-4-117}，
\[
C(L;\lambda_i)_
{\mathcal R^{\mathrm{sup}}_{l(\lambda_i)}}
\]
与
\[
C(L;\lambda_{i+1})_
{\mathcal R^{\mathrm{sup}}_{l(\lambda_i)}}
\]
作为滤过
$A_{n(\lambda_i),K(\lambda_i)}$ 代数
同伦等价。
注意
\[
C(L;\lambda_{i+1})_
{\mathcal R^{\mathrm{sup}}_{l(\lambda_i)}}
=
\frac{
C(L;\lambda_{i+1})_
{\mathcal R^{\mathrm{sup}}_{l(\lambda_{i+1})}}
}{
\mathfrak{IR}_{l(\lambda_i)}
},
\]
且
\[
C(L;\lambda_{i+1})_
{\mathcal R^{\mathrm{sup}}_{l(\lambda_{i+1})}}
\]
是一个
$A_{n(\lambda_{i+1}),K(\lambda_{i+1})}$ 代数。

因此，由适当的障碍理论，
\[
C(L;\lambda_i)_
{\mathcal R^{\mathrm{sup}}_{l(\lambda_{i+1})}}
\]
具有一个
$A_{n(\lambda_{i+1}),K(\lambda_{i+1})}$
代数结构，
它在
\[
C(L;\lambda_i)_
{\mathcal R^{\mathrm{sup}}_{l(\lambda_i)}}
\]
上诱导原有的
$A_{n(\lambda_i),K(\lambda_i)}$ 代数结构。
又有
\[
\varprojlim_{j\to\infty}
C(L;\lambda_i)_
{\mathcal R^{\mathrm{sup}}_{l(\lambda_j)}}
=
C(L;\lambda_i)_{\mathcal R^{\mathrm{sup}}}.
\]
故
$C(L;\lambda_i)_{\mathcal R^{\mathrm{sup}}}$
带有滤过 $A_\infty$ 代数结构。
于是，由引理~\ref{lem:7-4-116}
的 $A_\infty$ 版本\translatornote{源书误写“引理 7.4.100 的 $A_\infty$ 版本”；本书没有该引理，(7.4.100) 只是前一证明中的同余式。这里所需的正是从 $A$ 结构恢复 $L$ 同态的引理 7.4.116。见疑误 A-0582。}，
得到 $L_\infty$ 同态
\[
\mathfrak q^{\circ}:
C(M;\Lambda_{0,\mathrm{nov}})[1]
\longrightarrow
\operatorname{Der}B_G
\bigl(C(L;\Lambda_{0,\mathrm{nov}})[1]\bigr).
\]
这里置
\[
C(L;\Lambda_{0,\mathrm{nov}})
=C(L;\lambda_1)
\]
\translatornote{源书左端写成 $C(L;;\Lambda_{0,\mathrm{nov}})$，多一个分号；译稿依紧前同态及既定记号删去。见疑误 A-0583。}。
定理~\ref{thm:3-8-32} 得证。\hfill$\square$

注意，在上述证明中，
只在从 $\mathfrak q^{\mathrm{geo}}$
过渡到 $\mathfrak q$ 时
使用了定理~\ref{thm:7-4-41}。
下一小节将以更实质的方式使用该定理，
并通过分叉论证
（见第~\ref{subsec:7-2-14} 小节）
证明定理~\ref{thm:intro-Y}。

定理~\ref{thm:3-8-32} 的含幺版本
（它蕴含式 (3.8.35)）
可由前两个小节发展出的同调代数
以及第~\ref{sec:7-3} 节的方法推出。

使用定理~\ref{thm:3-8-32}、
推论~\ref{cor:7-4-40}
和定理~\ref{thm:7-4-41}，得到下述结论。

\begin{theorem}{7.4.118}\label{thm:7-4-118}
\liftIndex{I359}
存在严格、滤过且有隙的 $L_\infty$ 同态
\[
\mathfrak q^{\circ}:
H(M;\Lambda_{0,\mathrm{nov}})[1]
\longrightarrow
\operatorname{Der}B
\bigl(H(L;\Lambda_{0,\mathrm{nov}})[1]\bigr),
\]
使得把系数环约化为 $\mathbb Q$ 后，
$\mathfrak q^{\circ}$ 约化为由包含映射诱导的线性映射
\[
H(M;\mathbb Q)
\longrightarrow
\operatorname{Hom}\bigl(\mathbb Q,H(L;\mathbb Q)\bigr)
\cong H(L;\mathbb Q).
\]
也可把
$\operatorname{Der}B
\bigl(H(L;\Lambda_{0,\mathrm{nov}})[1]\bigr)$
替换为
$\operatorname{Der}B^{\mathrm{unit}}
\bigl(H(L;\Lambda_{0,\mathrm{nov}})[1]\bigr)$。
\end{theorem}

最后一个断言是式~\eqref{eq:3-8-34-1} 的推论。
